\section{Introduction}\label{sec:intro}

\paragraph{GNSS-denied navigation.} A small unmanned aerial vehicle (UAV) normally knows where it
is from satellite navigation (GNSS). When GNSS is jammed, spoofed or simply unavailable, it falls
back on its inertial navigation system (INS), which integrates measured velocity. Any error in
that velocity is integrated too: a constant velocity bias of only $0.3$\,m/s --- the default
INS grade in this study --- displaces the position estimate by $90$\,m after five minutes, and
the error keeps growing without bound.

\paragraph{Terrain-aided navigation.} Most aircraft carry two more sensors that together
measure something the INS cannot: a barometric altimeter gives the altitude above sea level and a
radar altimeter the height above the ground below. Their difference is the height of the terrain
directly under the aircraft. A digital elevation model (DEM) stored on board says what that height
should be at every candidate position, so a sequence of measured terrain heights --- a terrain
profile --- can be matched against the map to recover the position. This is terrain-aided
navigation (TAN)~\cite{bergman1999csm}. Early TAN systems used Kalman filters that linearise the
terrain around the current estimate~\cite{hostetler1983}; Bergman, Ljung and Gustafsson showed why
the problem is better treated as full Bayesian estimation~\cite{bergman1999csm,bergman1999thesis},
and particle filters became a standard tool for it~\cite{gustafsson2002,gustafsson2010}.

\paragraph{Why it is hard.} The measurement is a nonlinear function of position, and a strongly
non-unique one: one terrain height is matched by a whole contour line, and a short profile can
fit several places equally well. The posterior is therefore often \emph{multimodal}, and how
informative the measurements are depends entirely on the terrain --- rugged hills pin the position
down within metres, a flat plain hardly at all. A linearised Gaussian filter (the EKF) describes
the terrain by its slope at one point, which is meaningless when the position is uncertain by
hundreds of metres; a particle filter represents the full posterior and needs no linearisation,
at a computational cost.

\paragraph{This project.} A GPS-denied UAV is simulated flying over real Copernicus GLO-30
terrain~\cite{copernicus2022} of three contrasting regions of Sri Lanka --- the rugged central
highlands, the hills around Kandy and the flat dry-zone plains --- and its position is estimated by
a sequential importance resampling (SIR) particle filter, an extended Kalman filter (EKF) and INS
dead reckoning. The contributions are:
\begin{itemize}
  \item a vectorised, log-domain SIR particle filter with four resampling schemes, roughening,
        a heavy-tailed (Student-$t$) likelihood, an optional measurement-offset state and an
        online divergence monitor, plus a Rao-Blackwellised particle filter; each validated against
        the exact Kalman filter on planar terrain;
  \item an experiment framework with common random numbers, in which every filter and every
        configuration sees the same randomised flights and sensor noise, so differences are paired
        and tested (McNemar, bootstrap);
  \item seven experiments (\cref{sec:results}): the filters against the initial uncertainty on
        each terrain, terrain informativeness at the level of single flights, particle-filter
        design, robustness to model mismatch (outliers, map error, INS grade), consistency against
        the posterior Cramér--Rao bound, false-lock detection and recovery, and
        Rao-Blackwellisation;
  \item full reproducibility: study regions chosen by a rule fixed in advance, every number,
        table and figure generated by scripts from saved results, and the data traced by
        checksums.
\end{itemize}
\Cref{sec:method} derives the models and estimators, \cref{sec:impl} describes the
implementation and its validation, \cref{sec:results} presents the experiments,
\cref{sec:discussion} discusses what they show and their limitations, and \cref{sec:conclusion}
concludes.
